\part{Comment les nombres deviennent des octets}
\section{Binaire, hexadécimal et ordre des octets}
\subsection{Un octet n'a pas de type intrinsèque}
Un octet possède huit bits, donc 256 configurations. La suite
\texttt{01000001} vaut 65 en entier non signé et représente aussi
la lettre A en ASCII. Le type est une règle d'interprétation, pas
une propriété magique de ces huit bits. En hexadécimal, un chiffre
représente quatre bits : \texttt{41} désigne le même octet.

Pour des bits $b_i\in\{0,1\}$,
\begin{equation}
n=\sum_{i=0}^{k-1}b_i2^i.
\end{equation}
Le bit d'indice 0 est le bit de poids faible. Le nombre 291 vaut
\texttt{0x0123} et $1\times256+2\times16+3$.
Un affichage hexadécimal montre généralement une adresse à gauche,
les octets au centre et une interprétation texte à droite. Le texte
illisible ne signifie pas que le fichier est chiffré : des pixels
et des nombres binaires ne sont pas du texte.

\subsection{Little-endian et big-endian}
Pour l'entier 16 bits \texttt{0x1234}, big-endian écrit
\texttt{12 34}, little-endian écrit \texttt{34 12}.
L'ordre des octets n'est pas la même question que l'ordre des bits
dans un flux compacté. Un conteneur et son codec interne peuvent
avoir des conventions différentes.

Lire deux octets $a,b$ en big-endian donne $256a+b$ ; en
little-endian, $a+256b$. Interpréter \texttt{01 00} avec le mauvais
ordre transforme 1 en 256. Sur une image, cela produit souvent
une luminosité aberrante ou des motifs périodiques.

\subsection{Entiers signés et flottants}
En complément à deux, un entier signé de $b$ bits va de
$-2^{b-1}$ à $2^{b-1}-1$. Un résidu de compression peut être
négatif même si l'échantillon final est non signé.

Un flottant binaire normalisé peut être compris comme
\begin{equation}
(-1)^s\,2^{e-\mathrm{biais}}\left(1+\frac{m}{2^p}\right).
\end{equation}
Il distribue une précision approximativement relative sur une
large plage. Le type binary32 utilise 1 bit de signe, 8 d'exposant
et 23 de fraction ; binary16 utilise 1, 5 et 10. Il existe aussi
zéros, subnormaux, infinis et NaN. Les valeurs très proches de zéro
ne suivent pas toutes la formule normalisée ci-dessus.

Un flottant 16 bits n'est donc pas un entier 16 bits : il peut
représenter des valeurs négatives et supérieures à 1, mais ne
possède pas 65 536 niveaux uniformément répartis entre 0 et 1.
Dans un moteur, les NaN et infinis doivent être détectés : une
seule valeur non finie peut contaminer histogrammes et statistiques.

\section{Organisation d'une image en mémoire}
\subsection{Entrelacé et planaire}
Une image RGB peut être rangée
\texttt{R0 G0 B0 R1 G1 B1}, ou en trois plans
\texttt{R0 R1 ...}, \texttt{G0 G1 ...}, \texttt{B0 B1 ...}.
La première organisation facilite certaines opérations par pixel ;
la seconde facilite certaines opérations par canal et certains codecs.
RGB et BGR sont aussi deux ordres différents. Une permutation R/B
est un problème d'organisation, pas une balance des blancs.

Pour une image entrelacée de $C$ canaux, avec $s$ octets par composante :
\begin{equation}\label{eq:address}
\operatorname{adresse}(x,y,c)=A_0+y\,\mathrm{stride}+xCs+cs.
\end{equation}
Le \emph{stride}, ou pitch, est le nombre d'octets séparant les
débuts de deux lignes. Il peut dépasser $WCs$ pour respecter un
alignement ou parce que l'image est une vue dans un plus grand tampon.

\subsection{Exemple de lignes rembourrées}
Une ligne de 3 pixels RGB 8 bits contient 9 octets utiles.
Si chaque ligne est alignée sur 4 octets, le stride peut être 12.
Les 3 derniers octets sont du remplissage, pas un quatrième pixel.
Pour $H=2$, la mémoire allouée vaut alors 24 octets, pas 18.
Une boucle qui suppose des lignes contiguës de 9 octets lira mal
la deuxième ligne.

Une région d'intérêt peut conserver le stride de son parent.
Son pointeur commence plus loin, mais sa largeur n'explique pas
toute la distance entre ses lignes. En C++/OpenCV, ne présumer
la continuité d'un \texttt{cv::Mat} qu'après vérification ; dans Qt,
utiliser le pas réel de l'image. Une vue NumPy peut également être
non contiguë et même posséder un stride négatif.

\subsection{Origine, axes et orientation}
En traitement d'image, l'origine est souvent en haut à gauche,
$x$ croît vers la droite et $y$ vers le bas. En géométrie
mathématique, $y$ croît souvent vers le haut. Une rotation positive
peut donc sembler inversée si l'on confond transformation du repère
et transformation du contenu.

Il faut distinguer la grille stockée de l'orientation d'affichage.
Une photo portrait peut contenir une grille paysage avec une
instruction de rotation. Les huit orientations Exif couvrent
rotations et réflexions~\cite{exif}. Pour un exemple de rotation
horaire d'une grille $W\times H$ :
\[
x'=H-1-y,\qquad y'=x.
\]
Les dimensions deviennent $H\times W$. Si le logiciel tourne les
pixels tout en conservant une instruction de rotation non remise
à jour, le lecteur suivant peut effectuer une seconde rotation.

\subsection{Alpha droit et alpha prémultiplié}
L'alpha exprime une couverture ou opacité conventionnelle, pas une
quatrième couleur ni forcément une transparence physique complète.
Avec alpha droit, la couleur $C$ est stockée indépendamment de $a$.
Avec alpha prémultiplié, le tampon stocke $P=aC$.

Pour deux couches, avant-plan $f$ et arrière-plan $b$, en prémultiplié :
\begin{equation}
P_{\mathrm{out}}=P_f+(1-a_f)P_b,\qquad
a_{\mathrm{out}}=a_f+(1-a_f)a_b.
\end{equation}
Si l'on veut une couleur droite, diviser $P_{\mathrm{out}}$ par
$a_{\mathrm{out}}$ quand il est non nul. Les calculs de mélange
lumineux doivent se faire dans un espace linéaire approprié.
Dé-prémultiplier quand $a$ est presque nul amplifie les erreurs.
Des franges sombres autour des objets proviennent fréquemment
d'un mauvais contrat d'alpha ou d'une interpolation dans le
mauvais domaine.

\section{Coût mémoire et transfert GPU}
Pour $6000\times4000$ pixels :
\begin{center}
\begin{tabular}{lrr}
\toprule Représentation & Octets/pixel & Taille sans padding \\
\midrule Mosaïque stockée en uint16 & 2 & 48 MB \\
RGB uint16 & 6 & 144 MB \\
RGBA float32 & 16 & 384 MB \\
Deux tampons RGBA float32 & 32 & 768 MB \\
\bottomrule
\end{tabular}
\end{center}
Ici MB signifie $10^6$ octets ; 384 MB représentent environ
366,2 MiB, avec $1\,\mathrm{MiB}=2^{20}$ octets. Il faut ajouter
les intermédiaires, pyramides, masques, textures GPU et caches.
Le petit fichier JPEG de quelques mégaoctets peut ainsi produire
des centaines de mégaoctets en mémoire.

Un GPU ne lit pas nécessairement le TIFF ou le CR3 directement.
La chaîne usuelle comprend lecture, décompression CPU ou spécialisée,
allocation d'un tampon, transfert, traitement, puis éventuellement
retour CPU. Le coût total est
\[
T=T_{\text{lecture}}+T_{\text{décompression}}+
T_{\text{transferts}}+T_{\text{calcul}}+T_{\text{attentes}}.
\]
Un filtre très rapide ne compense pas toujours les transferts.
Mesurer chaque étage avant de conclure que « le GPU accélérera
l'ouverture ». Les formats de texture de rendu ne sont pas
nécessairement les formats photographiques de stockage.

\section{Un fichier est un conteneur structuré}
\subsection{Signature, en-tête, répertoires et charges utiles}
Une architecture générale peut être représentée ainsi :
\begin{center}
\begin{tikzpicture}[node distance=.45cm,
box/.style={draw=iprblue,rounded corners,text width=11cm,align=center,minimum height=.8cm}]
\node[box] (s) {Signature et version : quel contrat de lecture ?};
\node[box,below=of s] (h) {En-têtes et index : dimensions, types, offsets, compression};
\node[box,below=of h] (m) {Métadonnées et calibration : Exif, profils, orientation};
\node[box,below=of m] (d) {Charges utiles : pixels, coefficients, aperçus, masques};
\node[box,below=of d] (c) {Contrôles éventuels : longueurs, CRC, sommes de contrôle};
\draw[-{Latex}] (s)--(h);\draw[-{Latex}] (h)--(m);
\draw[-{Latex}] (m)--(d);\draw[-{Latex}] (d)--(c);
\end{tikzpicture}
\end{center}
Ce schéma est conceptuel : les sections peuvent être dispersées
dans le fichier, imbriquées ou répétées. Un TIFF utilise des
répertoires et offsets ; un PNG des chunks successifs ; un fichier
fondé sur ISO BMFF des boîtes imbriquées. Il n'existe pas un
en-tête universel de toutes les images.

\subsection{Codec, conteneur et extension}
Le conteneur organise les blocs ; le codec explique comment
reconstruire les échantillons de la charge utile. Une même
extension peut accepter plusieurs codecs. TIFF peut stocker des
valeurs non compressées ou différents flux compressés. CR3
désigne une famille de conteneurs Canon, non un simple tableau
RGB de convention fixe.

Renommer \texttt{photo.jpg} en \texttt{photo.dng} ne convertit rien.
Le type MIME, l'extension et la signature sont des indices à
croiser. Une signature valide prouve seulement que quelques
octets ressemblent à l'entrée d'un format, pas que tous ses blocs
sont corrects ou que son image se décodera.

\subsection{Offsets, longueurs et débordements}
Un offset est une distance en octets depuis une origine définie.
Cette origine peut être le début du fichier, celui d'un bloc ou
celui d'une structure interne. Une longueur peut compter les
données seules ou inclure l'en-tête. Ne pas deviner ces conventions.

Si un bloc commence à $o$ et contient $n$ octets dans un fichier
de taille $S$, vérifier $0\le o\le S$ puis $0\le n\le S-o$.
Cette seconde forme évite le débordement de $o+n$ dans un entier
de taille finie. De même, valider $W$, $H$, les canaux et le produit
avant d'allouer $WHCs$. Un fichier malveillant peut déclarer une
image gigantesque tout en ne contenant que quelques octets.

\subsection{CRC, hash et identité}
Un CRC détecte efficacement certaines corruptions accidentelles.
Un hash cryptographique comme SHA-256 sert à comparer des contenus
et à contrôler la provenance ; il ne prouve pas à lui seul
l'authenticité d'un auteur. Une signature cryptographique ajoute
une identité et une chaîne de confiance à vérifier.

Deux fichiers de hashes différents peuvent décoder les mêmes
pixels : ordre des métadonnées, compression et horodatages
diffèrent. À l'inverse, deux fichiers de même nom peuvent avoir
des contenus distincts. Il faut distinguer identité des octets,
égalité des pixels décodés et équivalence du rendu colorimétrique.

\resultat{Le fichier peut être compact, dispersé et compressé ; la
mémoire peut être large, alignée et flottante. Lire un fichier n'est
jamais simplement copier ses octets dans une texture.}
